% !PRINTED 354
\weberpage{354}

Da $a$ durch $\mathfrak{a}$ teilbar sein soll, so setzen wir
\[ a = \mathfrak{a}\mathfrak{m}' = \mathfrak{a}'\mathfrak{m} \]
und erhalten aus (1):
\weq{2}{\mathfrak{a} = m\mathfrak{m}.}
Also geht $m$ im Teiler des Ideals $\mathfrak{a}$ auf, und damit ist der Satz 3. bewiesen, da, wenn $\mathfrak{a}$ primär ist, sein Teiler $= 1$ sein muß.

4. \gesperrt{Satz. Ist $\mathfrak{m}$ ein primäres Ideal des Körpers $\Omega$, so ist jede ganze Zahl in $\Omega$ nach dem Modul $\mathfrak{m}$ mit einer rationalen Zahl kongruent.}

Nach Bd.\ II, \S\ 165 ist nämlich allgemein $N(\mathfrak{m})$ die Anzahl der nach $\mathfrak{m}$ inkongruenten Zahlen in $\Omega$. Ist nun $N(\mathfrak{m}) = a$ zugleich die kleinste durch $\mathfrak{m}$ teilbare \gesperrt{natürliche} Zahl, so sind die $a$ Zahlen
\[ 0, 1, 2, \ldots, a - 1 \]
alle inkongruent, und darunter sind alle Zahlklassen modulo $\mathfrak{m}$ vertreten.

Es sei jetzt $[Q]$ eine Ordnung mit dem Führer $Q$ und der Diskriminante $D$ und $(1, \theta)$ die Basis dieser Ordnung. Ferner sei $\mathfrak{m}$ ein \gesperrt{primäres} zu $Q$ teilerfremdes Ideal. Nach dem Satz 4. gibt es dann eine ganze rationale Zahl $t$, die der Bedingung
\[ \theta + t \equiv 0 \pmod{\mathfrak{m}} \]
genügt. Wir setzen
\weq{3}{\begin{aligned}
t &= \frac{b}{2}, && \text{wenn } D \equiv 0 \\
t &= \frac{1+b}{2}, && \text{„} \quad D \equiv 1
\end{aligned}
\qquad \text{(mod 4)}}
ist, und erhalten nach \S\ 96, (5):
\weq{4}{\theta + t = \frac{b + \sqrt{D}}{2} \equiv 0 \pmod{\mathfrak{m}},}
worin $b$ eine ganze rationale Zahl ist, die nach (3) der Bedingung
\[ b \equiv D \pmod{2} \]
genügt. Außerdem sei
\weq{5}{N(\mathfrak{m}) = a,}
also $a$ die kleinste durch $\mathfrak{m}$ teilbare natürliche Zahl.

Durch die Kongruenz (4) ist die Zahl $b$ nur nach dem Modul $2a$ bestimmt. Denn sind $b, b'$ zwei Zahlen, die der Bedingung (4) genügen, so ist $\frac{1}{2}(b' - b)$ durch $a$ teilbar.

Da ferner
\[ N\left(\frac{b + \sqrt{D}}{2}\right) = \frac{b^2 - D}{4} \]
%CONT>>
